\documentclass[10pt,a4paper]{amsart}
\usepackage[margin=19mm]{geometry}
\usepackage{amsmath,amssymb,mathtools}
\usepackage[scr=rsfs,cal=euler]{mathalfa}
\usepackage{xfrac}
\usepackage[unicode,hidelinks,bookmarks=false]{hyperref}
\newcommand{\ie}{i.e.\ }
\newcommand{\cf}{cf.\ }
\newcommand{\resp}{respectively\ }
\newcommand{\Subsection}{\subsection}
\providecommand{\nobreakdash}{\nobreak\hbox{-}}
\setlength{\emergencystretch}{2em}
\multlinegap0pt
\usepackage{bm}
\usepackage{bbm}
\usepackage{dsfont}
\usepackage{xspace}
\usepackage{cancel}
\usepackage{mathtools}
\usepackage{amsthm}
\makeatletter
\@ifundefined{prop}{\newtheorem{prop}{Proposition}}{}
\makeatother
\newcommand{\calB}{\mathcal B}
\title[Zeros of Stern polynomials in the complex plane]{Zeros of Stern polynomials in the complex plane}
\author{David Altizio}
\date{Source: arXiv:2511.03847v1 (CC BY 4.0)}
\begin{document}
\maketitle

\noindent\emph{Source and licence.} Zeros of Stern polynomials in the complex plane. Version arXiv:2511.03847v1 (CC BY 4.0), distributed under the Creative Commons Attribution 4.0 International licence, as recorded in the arXiv licence metadata for this exact version and shown on the abstract page https://arxiv.org/abs/2511.03847v1. Full rights notice: licenses/arxiv-2511.03847-CC-BY-4.0.txt.\par\smallskip

\noindent\textit{Editorial scope.} Counted unit: the Prop whose source label is \texttt{prop:ineq-smallcases} in the article (source0.tex lines 825--831, source label \texttt{prop:ineq-smallcases}), delivered together with the article's own complete original proof of it (source0.tex lines 833--873, one \texttt{proof} environment of 41 lines). No mathematics is written, repaired or re-derived: statement and proof are the article's own text line for line. Required context retained uncounted: none. Every definition, display and label of those lines is retained unchanged. Omitted from this package and not redistributed: 1--824, 832--832, 874--1127. Nothing inside a retained excerpt is omitted or altered: every retained body is delivered exactly as the article's lines are written, so no omission receipt of any kind accompanies this package. Citations: the retained excerpts carry no citation command at all, so no bibliography entry of the article is transcribed and no reference is redistributed. Reference adaptation: none was needed and none was made; every reference inside the delivered unit resolves inside it, so no unresolved reference or citation is produced. Printed slips kept literal: no printed word, symbol, hypothesis or number of the counted statement or of its proof was altered. Layout and class: the article's own class is \texttt{article} with packages including multicol, amsmath, amssymb, amsthm, array, booktabs, mathtools, tabularray, mathpazo, geometry, caption, chngcntr; the wrapper is the lane's pinned \texttt{amsart} editorial wrapper at 10pt on A4, which restates the theorem-like environments the retained text needs and adds the article's own macro definitions that the retained text needs (\texttt{\textbackslash calB}), with extra wrapper packages (bm, bbm, dsfont, xspace, cancel, mathtools). The delivered numbering is the unit's own. Assets: no retained span contains a figure environment, a graphics-inclusion command, a PDF-page inclusion, a table or a TikZ picture; no image file is copied into this package, the package declares no asset dependency and no image file is redistributed with it. Licence: version arXiv:2511.03847v1 of this article is distributed under the Creative Commons Attribution 4.0 International licence, as recorded in the arXiv licence metadata for this exact version and shown on the abstract page https://arxiv.org/abs/2511.03847v1. A full rights notice is delivered at licenses/arxiv-2511.03847-CC-BY-4.0.txt. Characters: every retained excerpt is pure ASCII, so the delivered engine, XeLaTeX with its default Latin Modern text font, reports no missing character.
\begin{prop}\label{prop:ineq-smallcases}
	Suppose $z\in \calB$.  Then
	\begin{align}
		|(z)_5| &= |1+z+z^2+z^3+z^4| \geq 5|z|,\label{eqn:min-5}\\
		|(z)_6| &= |1+z+z^2+z^3+z^4+z^5| \geq 6|z|.\label{eqn:min-6}
	\end{align}
\end{prop}

\begin{proof}
	We first prove \eqref{eqn:min-5}.  Observe that the function $f(z) = \tfrac{z}{(z)_5}$ is analytic in $\calB$, so by the Maximum Modulus Principle we may assume $z\in\partial B$, i.e. $|z-2| = 1$.  Write $z = re^{i\theta}$.  Rewrite the given inequality to $|z^5 - 1| \geq 5|z| |z-1|$ and square both sides; this yields the equivalent inequality
	\begin{equation}\label{eqn:polar-coord}
		r^{10} - 2r^5\cos(5\theta) + 1 \geq 25r^2(r^2 - 2r\cos\theta + 1).
	\end{equation}
	Observe that $|z-2| = 1$ implies $r^2 - 4r\cos\theta + 4 = 1$, or $\cos\theta = \frac{r^2+3}{4r}$.  Under this substitution, the right hand side of \eqref{eqn:polar-coord} becomes
	\[
	25r^2\left[r^2 - 2r\cdot \frac{r^2+3}{4r} + 1\right] = \frac{25}2(r^4 - r^2).
	\]
	Analogously, the left hand side of \eqref{eqn:polar-coord} becomes
	\begin{align*}
		r^{10} - 2r^5\cos(5\theta) + 1 &= r^{10} - 2r^5\left(16\cos^5\theta - 20\cos^3\theta + 5\cos\theta\right) + 1\\
		&= r^{10} - 2r^5\left[16\left(\frac{r^2+3}{4r}\right)^5 - 20\left(\frac{r^2+3}{4r}\right)^3 + 5\left(\frac{r^2+3}{4r}\right)\right] + 1\\
		&= r^{10} - 32\left(\frac{r^2+3}4\right)^5 - 20r^2\left(\frac{r^2+3}{4}\right)^3 + 5r^4\left(\frac{r^2+3}4\right)+1\\
		&= \frac1{32}(31r^{10} + 5r^8 + 10r^6 + 30r^4 + 135r^2 - 211).
	\end{align*}
	
	It follows that \eqref{eqn:polar-coord} is equivalent to the inequality 
	\begin{equation}\label{eqn:reduction-inequality}
		\frac1{32}(31r^{10} + 5r^8 + 10r^6 + 30r^4 + 135r^2 - 211) \geq \frac{25}2(r^4 - r^2).
	\end{equation}
	But \eqref{eqn:reduction-inequality} follows from the unexpected factorization
	\[
	\text{LHS} - \text{RHS} = \frac{(r^2-1)^3(31r^4 + 98r^2 + 211)}{32} \geq 0.
	\]
	This proves \eqref{eqn:min-5}.
	
	The proof of \eqref{eqn:min-6} is similar, so we only sketch the details.  By the Maximum Modulus Principle we may assume $|z-2| = 1$.  Squaring both sides yields the equivalent inequality
	\begin{equation}\label{eqn:polar-coord-2}
		r^{12} - 2r^6\cos(6\theta) + 1 \geq 36r^2(r^2 - 2r\cos\theta + 1).
	\end{equation}
	Substituting $\cos \theta = \tfrac{r^2+3}{4r}$ and expanding yields the inequality
	\[
	\frac{1}{64} \left(63 r^{12}+6 r^{10}+9 r^8+20 r^6+81 r^4+486 r^2-665\right) \geq 18(r^4 - r^2).
	\] 
	Here, we have the factorization
	\begin{equation}\label{eqn:ineq-2}
		\text{LHS} - \text{RHS} =  \frac{(r^2-1)^2(63r^8+132r^6+210r^4+308r^2-665)}{64}.
	\end{equation}
	The polynomial $g(x)\coloneqq  63x^8+132x^6+210x^4+308x^2-665$ is increasing on $[0,\infty)$ and $g(1) = 48$.  Because $r\in[1,3]$, we know $g(r) \geq 0$, and thus the right hand side of \eqref{eqn:ineq-2} is always nonnegative.
\end{proof}

\end{document}
